\subsection{\texorpdfstring{第~\ref{sec:hamiltonian-mechanics}~节}{Hamilton 力学}}\label{subsec:solutions-hamiltonian-mechanics}

\begin{description}[font=\normalfont,leftmargin=2em,labelwidth=1.5em]
\RGsolution{ex:verify-canonical-symplectic-form}{verify-canonical-symplectic-form}
考虑 $T^*M$ 上的典范辛形式, 它在局部坐标 $(x^1,\ldots,x^n,p_1,\ldots,p_n)$ 下由
\[
  \omega=\sum_{j=1}^n\mathrm{d}p_j\wedge\mathrm{d}x^j
\]
给出. 这个形式显然是闭的. 为证明它在 $\alpha\in T^*M$ 处非退化, 取 $v\in T_\alpha(T^*M)$, 使得对每个 $w\in T_\alpha(T^*M)$ 有 $\omega(v,w)=0$. 写成
\[
  v=\sum_{j=1}^n a^j\frac{\partial}{\partial x^j}+b_j\frac{\partial}{\partial p_j},
\]
并对 $i\in\{1,\ldots,n\}$ 取 $w:=\frac{\partial}{\partial p_i}$, 得
\[
  0=\omega(v,w)=\sum_{j=1}^n(\mathrm{d}p_j\otimes\mathrm{d}x^j-\mathrm{d}x^j\otimes\mathrm{d}p_j)
     \left(v,\frac{\partial}{\partial p_i}\right)=-a^i.
\]
若改取 $w:=\frac{\partial}{\partial x^i}$, 其中 $i\in\{1,\ldots,n\}$, 则
\[
  0=\omega(v,w)=\sum_{j=1}^n(\mathrm{d}p_j\otimes\mathrm{d}x^j-\mathrm{d}x^j\otimes\mathrm{d}p_j)
     \left(v,\frac{\partial}{\partial x^i}\right)=b_i.
\]
因此 $v=0$, 所以 $\omega$ 非退化. 最后, 形式
\begin{align*}
  \omega^n&=\omega\wedge\cdots\wedge\omega
     =\left(\sum_{j=1}^n\mathrm{d}p_j\wedge\mathrm{d}x^j\right)\wedge\cdots\wedge
       \left(\sum_{j=1}^n\mathrm{d}p_j\wedge\mathrm{d}x^j\right)\\
  &=n!\,\mathrm{d}p_1\wedge\mathrm{d}x^1\wedge\cdots\wedge\mathrm{d}p_n\wedge\mathrm{d}x^n\\
  &=n!(-1)^{\frac{n(n-1)}2}
        \mathrm{d}p_1\wedge\cdots\wedge\mathrm{d}p_n\wedge\mathrm{d}x^1\wedge\cdots\wedge\mathrm{d}x^n
        \ne0
\end{align*}
是 $T^*M$ 上的一个体积形式.

\RGsolution{ex:conservative-system-with-magnetic-term}{conservative-system-with-magnetic-term}
\begin{enumerate}[label=(\alph*),ref=\alph*,leftmargin=*]
\item 在局部坐标中, 有
\begin{align*}
  L(x^1,\ldots,x^n,v^1,\ldots,v^n)
  &=\frac12\sum_{i,j=1}^n g_{ij}(x^1,\ldots,x^n)v^iv^j\\
  &\quad+\sum_{i=1}^n\alpha_i(x^1,\ldots,x^n)v^i-U(x^1,\ldots,x^n)
     =K+C-U,
\end{align*}
其中 $\alpha=\sum_{i=1}^n\alpha_i\,\mathrm{d}x^i$ 且 $C=\sum_{i=1}^n\alpha_i v^i$. 我们已知
\[
  \mu\left(\frac{D\dot c}{\mathrm{d}t}(t)\right)+(\mathrm{d}U)_{c(t)}
\]
在局部坐标中由下式给出:
\begin{align*}
  &\sum_{i=1}^n\left[\frac{\mathrm{d}}{\mathrm{d}t}
    \left(\frac{\partial K}{\partial v^i}(x(t),\dot{x}(t))\right)
    -\frac{\partial K}{\partial x^i}(x(t),\dot{x}(t))\right]\mathrm{d}x^i\\
  &\qquad+\sum_{i=1}^n\frac{\partial U}{\partial x^i}(x(t))\,\mathrm{d}x^i.
\end{align*}
此外,
\begin{align*}
  &\sum_{i=1}^n\left[\frac{\mathrm{d}}{\mathrm{d}t}
    \left(\frac{\partial C}{\partial v^i}(x(t),\dot{x}(t))\right)
    -\frac{\partial C}{\partial x^i}(x(t),\dot{x}(t))\right]\mathrm{d}x^i\\
  &\quad=\sum_{i=1}^n\left[\frac{\mathrm{d}}{\mathrm{d}t}(\alpha_i(x(t)))
        -\sum_{j=1}^n\frac{\partial\alpha_j}{\partial x^i}(x(t))\dot{x}^j(t)\right]\mathrm{d}x^i\\
  &\quad=\sum_{i=1}^n\left[\sum_{j=1}^n\left(\frac{\partial\alpha_i}{\partial x^j}(x(t))
      -\frac{\partial\alpha_j}{\partial x^i}(x(t))\right)\dot{x}^j(t)\right]\mathrm{d}x^i,
\end{align*}
因此 $\iota(\dot c(t))\mathrm{d}\alpha$ 在局部坐标中由
\[
  \sum_{i=1}^n\left[\frac{\mathrm{d}}{\mathrm{d}t}\left(\frac{\partial C}{\partial v^i}(x(t),\dot{x}(t))\right)
    -\frac{\partial C}{\partial x^i}(x(t),\dot{x}(t))\right]\mathrm{d}x^i
\]
给出. 于是 Euler–Lagrange 方程等价于
\begin{align*}
  &\sum_{i=1}^n\left[\frac{\mathrm{d}}{\mathrm{d}t}
      \left(\frac{\partial L}{\partial v^i}(x(t),\dot{x}(t))\right)
        -\frac{\partial L}{\partial x^i}(x(t),\dot{x}(t))\right]\mathrm{d}x^i=0\\
  &\iff \sum_{i=1}^n\left[\frac{\mathrm{d}}{\mathrm{d}t}
      \left(\frac{\partial(K+C)}{\partial v^i}(x(t),\dot{x}(t))\right)\right.\\
  &\qquad\left.-\frac{\partial(K+C)}{\partial x^i}(x(t),\dot{x}(t))
                   +\frac{\partial U}{\partial x^i}(x(t))\right]\mathrm{d}x^i=0\\
  &\iff\mu\left(\frac{D\dot c}{\mathrm{d}t}(t)\right)+\iota(\dot c(t))+\mathrm{d}\alpha(\mathrm{d}U)_{c(t)}=0,
\end{align*}
所求结论随即成立.
\item 有
\begin{align*}
  \frac{\mathrm{d}E}{\mathrm{d}t}(t)
  &=\frac{\mathrm{d}}{\mathrm{d}t}\left(\frac12\langle\dot c(t),\dot c(t)\rangle+U(c(t))\right)\\
  &=\left\langle\frac{D\dot c}{\mathrm{d}t}(t),\dot c(t)\right\rangle+(\mathrm{d}U)_{c(t)}(\dot c(t))\\
  &=\mu\left(\frac{D\dot c}{\mathrm{d}t}\right)(\dot c)+\mathrm{d}U(\dot c)
       =-\iota(\dot c)\mathrm{d}\alpha(\dot c)=-\mathrm{d}\alpha(\dot c,\dot c)=0.
\end{align*}
\item 在局部坐标中, Legendre 变换由
\[
  p_i=\frac{\partial L}{\partial v^i}=\sum_{j=1}^n g_{ij}v^j+\alpha_i
  \quad(i=1,\ldots,n)
\]
给出, 且容易求逆:
\[
  v^i=\sum_{j=1}^n g^{ij}(p_j-\alpha_j)\quad(i=1,\ldots,n).
\]
这说明 Lagrange 函数是超正则的.
\item 作为切丛上的函数, Hamilton 函数为
\begin{align*}
  H&=\sum_{i=1}^n v^i\frac{\partial L}{\partial v^i}-L
    =\sum_{i,j=1}^n g_{ij}v^iv^j+\sum_{i=1}^n\alpha_iv^i-L\\
    &=\frac12\sum_{i,j=1}^n g_{ij}v^iv^j+U.
\end{align*}
因此, 作为余切丛上的函数, 它由
\begin{align*}
  H&=\frac12\sum_{i,j,k,l=1}^n g_{ij}g^{ik}g^{jl}(p_k-\alpha_k)(p_l-\alpha_l)+U\\
   &=\frac12\sum_{i,j=1}^n g^{ij}(p_i-\alpha_i)(p_j-\alpha_j)+U
\end{align*}
给出, 而 Hamilton 方程为
\[
  \left\{\begin{aligned}
    \dot{x}^i
      &=\sum_{j=1}^n g^{ij}(p_j-\alpha_j),\\
    \dot p_i&=-\frac12\sum_{k,l=1}^n\frac{\partial g^{kl}}{\partial x^i}
          (p_k-\alpha_k)(p_l-\alpha_l)\\
      &\quad+\sum_{k,l=1}^n g^{kl}\frac{\partial\alpha_k}{\partial x^i}(p_l-\alpha_l)
           -\frac{\partial U}{\partial x^i}.
  \end{aligned}\right.
\]
\end{enumerate}

\RGsolution{ex:density-of-poincare-recurrent-points}{density-of-poincare-recurrent-points}
若满足对某个 $t\ge T$ 有 $\psi_t(p)\in U$ 的点 $p\in U$ 所成的集合不在 $U$ 中稠密, 则存在开集 $V\subset U$, 使得对所有 $t\ge T$, 都有 $\psi_t(V)\cap U=\varnothing$. 但将 Poincaré 回归定理应用于 $V$ 可保证存在一点 $p\in V$, 使得对某个 $t\ge T$ 有 $\psi_t(p)\in V\subset U$, 因而 $p\in\psi_t(V)\cap U$.

\RGsolution{ex:recurrence-of-unit-speed-geodesics}{recurrence-of-unit-speed-geodesics}
测地流对应于局部坐标中由
\[
  H(x^1,\ldots,x^n,p_1,\ldots,p_n)=\frac12\sum_{i,j=1}^n g^{ij}(x^1,\ldots,x^n)p_ip_j
\]
给出的 Hamilton 函数. 容易看出, 对流给出的每条测地线 $c:\mathbb{R}\to M$, $H$ 的 (守恒) 值就是 $\frac12\lVert\dot c(t)\rVert^2$. 由于 $M$ 紧, 集合 $K=H^{-1}([\frac18,\frac12])$ 也是紧的, 因为它与
\[
  \left\{v\in TM\ \middle|\ \frac12\le\lVert v\rVert\le1\right\}
\]
微分同胚. 考虑开集
\[
  U=\left\{\alpha\in T^*M\ \middle|\ \frac18<H(\alpha)<\frac12
                                     \text{ 且 }\pi(\alpha)\in B\right\}\subset K.
\]
由 Poincaré 回归定理, 存在 $\alpha\in U$, 使得对某个 $t_1\ge T$ 有 $\psi_{t_1}(\alpha)\in U$. 现在 $\psi_t(\alpha)$ 在 $M$ 上的投影是测地线 $\gamma(t)$, 且 $\gamma(0)\in B$ 和 $\gamma(t_1)\in B$. 此外, 对某个 $\frac12<k<1$, 有 $\lVert\dot\gamma(t)\rVert=k$. 因此 $c(t):=\gamma(t/k)$ 是一条满足 $\lVert\dot c(t)\rVert=1$ 的测地线, 它满足 $c(0)\in B$ 和 $c(t_1/k)\in B$, 即对某个 $t>T$ 有 $c(t)\in B$.
\end{description}
